\chapter{Sign information}
\label{chap:signs}

We record information on the sign of~$c_n$ and of $c_{n+1} + c_n$.

\begin{lemma}[Signs]
  \label{lem:signs}
  \uses{def:cseq, def:chars, def:phin, def:phi0tau}
  \lean{QuadraticIterates.cSeq_pos_of_pos, QuadraticIterates.cSeq_neg_of_neg_one_lt,
    QuadraticIterates.neg_le_cSeq_of_le_neg_two, QuadraticIterates.le_cSeq,
    QuadraticIterates.cSeq_le_sq_add, QuadraticIterates.hasCharPattern_mulSingle_one_of_pos,
    QuadraticIterates.hasCharPattern_neg_mulSingle_one_of_neg_one_lt}
  \leanok
  \begin{enumerate}
    \item If $a > 0$, then $c_1 = -a < 0$ and $c_n > 0$ for all $n \ge 2$. In particular, $a$ has
          $\phi_1 = \sgn$.
    \item If $-1 < a < 0$, then $c_1 = -a > 0$ and $c_n < 0$ for all $n \ge 2$. In particular,
          $a$ has $\tau = \sgn$.
    \item If $a \le -2$, then $c_n \ge -a > 0$ for all $n \ge 1$.
  \end{enumerate}
  \formalnote{The signs in a linearly ordered commutative ring, the patterns of~$\sgn$ for any
  $a \in \QQ$; (2) follows from \lname{cSeq\_le\_sq\_add}.}
\end{lemma}
\begin{proof}
  \leanok
  (1) $c_{n+1} = c_n^2 + a > 0$ for all $n \ge 1$. (2) $c_2 = a(a + 1) < 0$; if $c_n < 0$, then
  $c_{n+1} = c_n^2 + a \le a^2 + a < 0$, since $c_n = c_{n-1}^2 + a \ge a$ gives $c_n^2 \le a^2$.
  (3) $c_1 = -a > 0$; if $c_n \ge -a$, then $c_{n+1} = c_n^2 + a \ge a^2 + a \ge -a$.
\end{proof}

\begin{definition}[The number~$\alpha$]
  \label{def:alpha}
  \leanok
  We write~$\alpha$ for the unique real root of $x^3 + 2x^2 + 2x + 2$; note that
  $\alpha \approx -1.543689$ and that $a > \alpha$ is equivalent to $a^3 + 2a^2 + 2a + 2 > 0$.
  \formalnote{No definition: the statements use the inequality $a^3 + 2a^2 + 2a + 2 > 0$, for
  $a = r/s$ in the form $r^3 + 2r^2 s + 2rs^2 + 2s^3 > 0$.}
\end{definition}

\begin{lemma}[Sums of successive terms]
  \label{lem:sumsign}
  \uses{def:cseq, def:alpha, lem:signs}
  \lean{QuadraticIterates.add_cSeq_two_neg, QuadraticIterates.cSeq_add_cSeq_succ_neg,
    QuadraticIterates.cSeq_add_cSeq_succ_neg_of_lt, quadratic_le_max_of_nonneg_of_mem_Icc}
  \leanok
  Assume that $a \in (\alpha, 0)$. Then $c_{n+1} + c_n < 0$ for all $n \ge 2$.
\end{lemma}
\begin{proof}
  \leanok
  For $-2 \le a \le 0$, $c_3 \le c_2$ (as $c_3 - c_2 = a^3(a + 2)$), so the interval
  $[a, a^2 + a]$ is stable under~$f$: as~$f$ is convex from below, it takes its maximum on an
  interval at an endpoint, here~$c_2$. So $a \le c_n \le c_2$ for $n \ge 2$. The conditions
  $a^3 + 2a^2 + 2a + 2 > 0$ and $a < 0$ give $c_2^2 + c_2 + a < 0$. If $c_n \ge -1/2$, then
  $c_{n+1} + c_n = c_n^2 + c_n + a \le c_2^2 + c_2 + a < 0$; otherwise
  $c_{n+1} + c_n \le a^2 + a + a = a(a + 2) < 0$.
\end{proof}

For $a = \alpha$, $c_3 + c_2 = 0$, so the condition in the lemma is sharp. There is another
interval where the sign of~$c_n$ can be controlled in a useful way.

\begin{definition}[The number~$\alpha'$]
  \label{def:alphaprime}
  \leanok
  We write~$\alpha'$ for the larger real root of $x^6 + 3x^5 + 3x^4 + 3x^3 + 2x^2 + 1$. We have
  $\alpha' \approx -1.310703$, and $\alpha' < a < -1$ is equivalent to $c_3 < 0 < c_4$. ($c_4$
  is $a(a + 1)$ times the polynomial evaluated at~$a$.)
  \formalnote{No definition: the statements use $a < -1$, $c_3 < 0$ and $0 < c_4$, for $a = r/s$
  in the form $r < -s$ and $r w_3 < 0 < r w_4$.}
\end{definition}

\begin{lemma}[A sign window below~$-1$]
  \label{lem:window}
  \uses{def:cseq, def:alphaprime}
  \lean{QuadraticIterates.cSeq_mem_Icc_of_even, QuadraticIterates.cSeq_mem_Icc_of_odd,
    QuadraticIterates.cSeq_three_lt_neg_one}
  \leanok
  Assume that $\alpha' < a < -1$. Then for all $k \ge 1$: $0 < c_4 \le c_{2k} \le c_2$ and
  $a \le c_{2k+1} \le c_3 < -1$. So $\bigl(\sgn(c_n)\bigr)_{n \ge 1} = (1, 1, -1, 1, -1, 1,
  \dots)$.
  \formalnote{With the hypotheses \lname{a < -1}, \lname{cSeq a 3 < 0}, \lname{0 < cSeq a 4}.}
\end{lemma}
\begin{proof}
  \leanok
  \uses{lem:signs}
  $c_4 = c_3^2 + a > 0$ gives $c_3^2 > -a > 1$, so $c_3 < -1$; and $c_2 = a(a+1) > 0$. On
  $B_1 = [a, c_3] \subset (-\infty, 0)$ the map~$f$ is decreasing, so $f(B_1) = [c_4, c_2]$, using
  $f(a) = c_2$; on $B_2 = [c_4, c_2] \subset (0, \infty)$ it is increasing, so
  $f(B_2) = [c_5, c_3] \subseteq B_1$, as $c_5 \ge a$. Also $c_2 \in B_2$ (as $c_3 \in [a, 0)$
  gives $c_4 \le a^2 + a$) and $c_3 \in B_1$; induction on~$k$.
\end{proof}

So under the assumption of the lemma, $\sgn$ can serve as the character~$\psi$ in
Lemma~\ref{lem:jump2} when $n \ge 2$.
